\section{Local Carry Laws}

Carry normalization is governed by a local threshold rule. At degree \(k\), let
\[
    s_k = c_k + r_k,
\]
where \(c_k\) is the raw coefficient and \(r_k\) is the incoming carry. Then
\[
    d_k = s_k \bmod B,\qquad
    r_{k+1} = \left\lfloor \frac{s_k}{B} \right\rfloor.
\]

\begin{proposition}[Local Carry Activation]
The outgoing carry at degree \(k\) is active if and only if
\[
    c_k + r_k \ge B.
\]
\end{proposition}

\begin{proof}
By definition,
\[
    r_{k+1} = \left\lfloor \frac{c_k+r_k}{B} \right\rfloor.
\]
Since \(c_k+r_k\) is a nonnegative integer and \(B\ge 2\), the floor is
positive exactly when \(c_k+r_k \ge B\). Thus \(r_{k+1}>0\) if and only if
\(c_k+r_k \ge B\).
\end{proof}

\paragraph{Carry cascade length.}
A carry cascade is a maximal consecutive block of active carry edges
\[
    k \to k+1.
\]
The cascade length is the number of active edges in such a block. The maximum
carry cascade length records the longest uninterrupted carry propagation event.

\paragraph{Value-weighted carry flow.}
The value-weighted carry flow is
\[
    \operatorname{VWC}(c) = \sum_k r_{k+1} B^{k+1}.
\]
It measures transported carry mass in positional value units rather than raw
edge-count units. A normalized version divides this quantity by
\(\max(\sum_k c_k B^k,1)\).

Carry is local in recurrence but global in consequence because outgoing carry
changes the next degree. Thus \(c_k\) alone does not determine activation:
the correct local state variable is \(s_k=c_k+r_k\).
